\documentclass[a4paper,oneside, margin=1.4in]{book}

\usepackage{lipsum}
\usepackage[hidelinks]{hyperref}
\usepackage{titletoc}
\usepackage{amsmath}
\usepackage{geometry}
\usepackage{graphicx}
\graphicspath{ {.} }
\geometry{a4paper, margin=1in}
\titlecontents*{chapter}
  [0pt]% <left>
  {}
  {\chaptername\ \thecontentslabel\quad}
  {}
  {\bfseries\hfill\contentspage}    

\usepackage{bookmark}
\usepackage{etoolbox}

\makeatletter
\newcommand*{\AddChapterPrefixInBookmarks}{%
  \if@mainmatter
    \ifnum\bookmarkget{level}=0 %
      \preto\bookmark@text{\@chapapp\space}%
    \fi
  \fi
}
\makeatother

\bookmarksetup{
  numbered,
  addtohook=\AddChapterPrefixInBookmarks,
}

% Workaround for numbered sections in unnumbered
% chapter "Introduction" to avoid chapter number
% zero.
\renewcommand*{\thesection}{%
  \ifcase\value{chapter}%
  \else
    \thechapter.%
  \fi
  \arabic{section}%
}

\title{My document}

\begin{document}
\frontmatter

\begin{titlepage}
    \begin{center}
        \vspace*{1cm}
            
        \Huge
        \textbf{Statistics Question Bank}
            
        \vspace{0.5cm}
        \huge
        Second Paper
            
        \vspace{1.5cm}
            
        \textbf{Abdullah Al Mahmud}

     \vspace{1.5cm}

	\Large 
	Updated on: \today
            
        \vfill
            

            
        \vspace{0.8cm}
            
\includegraphics[width=1cm]{logo}
            
        \Large
        www.statmania.info\\
            
    \end{center}
\end{titlepage}


\tableofcontents


\mainmatter
\chapter{Probability} 
\section{Creative Questions}

\begin{enumerate}
  \item
  \textbf{It is observed that in a college, there are 100 students, of whom 30 play football, 40 play cricket, and 20 play both.}
 
  \begin{enumerate}
    \item
	What is the range of probability? \hfill 1
    \item
    	What is the relationship between independence and mutual excluvity?  \hfill 2
    \item
    	Are the probabilities of playing cricket and that of football independent? Prove. \hfill 3
     \item
     	If a student is selected randomly, and if he does not play cricket, what is the probability that \\ he plays football? \hfill 4
  \end{enumerate}

\item
	  \textbf{Sadman has an urn with 5 red and 4 white balls. He has randomly drawn two balls from the urn.} 
  
  \begin{enumerate}
    \item
	What is the probability of an uncertain event? \hfill 1
    \item
	Write the third axiom of probability. \hfill 2
    \item  
	What is the probability that both the balls drawn by Sadman are white? \hfill 3
    \item
	Are the probabilities of both balls being same color and different color equal? Analyze. \hfill 4
  \end{enumerate}

 \item
	  \textbf{Two dice are thrown together. The dice are named A and B.} 
  
  \begin{enumerate}
    \item
	What is P(A=7)? \hfill 1
    \item
	Create the sample space. \hfill 2
    \item  
	What is the probability that the outcomes of A \& B are different? \hfill 3
    \item
	Determine the probability that the summation of outcome of two dice is a prime number. \hfill 4
  \end{enumerate}

 \item
	  \textbf{A magician draws two cards from a pack (i) with replacement and then (ii) without replacement. The cards were well-shuffled before drawing.} 
  
  \begin{enumerate}
    \item
	What is the probability of an impossible event? \hfill 1
    \item
	How to determine the probability of a joint event?  \hfill 2
    \item  
	As per (i), what is the probability that the cards have different color? \hfill 3
    \item
	As per (ii), what is the probability that the cardsare aces of same color?  \hfill 4
  \end{enumerate}

 \item
	  \textbf{$P(A) = \frac{3}{10}, P(B) = \frac 25, P(B\cup A) = \frac12$} 
  
  \begin{enumerate}
    \item
	What is an independent event? \hfill 1
    \item
	What is the relationship between independency and mutual excluvity? \hfill 2
    \item  
	Find $P(A \vert B)$ and $P(B \vert A)$ \hfill 3
    \item
	Verify the equality mathematically \& empirically: $P(B) = P(A) \cdot P(B \vert A) + P(\bar A) \cdot P(B \vert \bar A)$ \hfill 4
  \end{enumerate}

 \item
  \textbf{Sakib has recently graduated from the University of Dhaka. he applies to two firms - EduCube \& Digic- for a Data Analyst job. The probability of hiring by EduCube is 0.8 and by Digic is 0.4. The probability that none hires is 0.5.} 
  
  \begin{enumerate}
    \item
	What is a sample space? \hfill 1
    \item
	Explain how to find $P(\bar A \cap B)$ using Venn Diagram. \hfill 2
    \item  
	Find the probability of hiirng by by Digic but not by EduCube. \hfill 3
    \item
	Find the probability that no firm will reject him. \hfill 4
  \end{enumerate}

 \item
	  \textbf{Recently there is an increase in the number of electronic medias in Bangladesh. A professor stated in the class room that very few people now resort to print media for news. A research indicates 70\% people collect news from electronic media, 60\% from print media, and 50\% from both.} 
  
  \begin{enumerate}
    \item
	What is an impossible event? \hfill 1
    \item
	Write the event "None of the two occurs" in two different notations. \hfill 2
    \item  
	What is the probability of getting news from at most one type of media? \hfill 3
    \item
	Is the professor correct in his/her statement? Analyze. \hfill 4
  \end{enumerate}
  
  \item
	  \textbf{A coin is tossed five times. The number of heads appearing from the tosses is considered a discrete random random variable.} 
  
  \begin{enumerate}
    \item
	What is a discrete random variable? \hfill 1
    \item
	Are probability distributions and frequency distributions similar? Show with an example. \hfill 2
    \item  
	Find the probability distribution from the stem and show in a table.\hfill 3
    \item
	Find the probability that a head will appear in more than 3 tosses.  \hfill 4
  \end{enumerate}

  
   \item
	  \textbf{Two dices-- one red and one blue-- are flung together. The dice are absolutely neutral and independent.} 
  
  \begin{enumerate}
    \item
	What is a simple event? \hfill 1
    \item
	Question \hfill 2
    \item  
	Question \hfill 3
    \item
	Question \hfill 4
  \end{enumerate}

  \end{enumerate}

\section{Short Questions}

  \begin{enumerate}
    \item
	Question \hfill 1
    \item
	Question \hfill 2
    \item  
	Question \hfill 3
    \item
	Question \hfill 4
  \end{enumerate}

\chapter{Random Variable and Probability Function} 
\section{Creative Questions}

  \begin{enumerate}

 \item
	  \textbf{The probability distribution of a discrete random variable X is given below:} 

	  \begin{table}[h]
	  \begin{center}
\begin{tabular}{lllllll}
x    & -2   & -1 & 0   & 1 & 3 & 4   \\ \hline
P(x) & 0.1 & k & 2k & 3k & 4k & 0.2
\end{tabular} 
\end{center}	
\end{table}
  
  \begin{enumerate}
    \item
	What is $\Sigma P(x)$? \hfill 1
    \item
	Find the value of k. \hfill 2
    \item  
	Find $P(X \geq 0) \& P(X < 1)$ \hfill 3
    \item
	Find the cumulative distribution function, F(X) and F(2) and explain. \hfill 4
  \end{enumerate}

 \item
	  \textbf{The joint probability function of two random variables X \& Y is given below:}

$P(x,y) = \frac{1}{21}(x+y); x=1,2,3$ \& $y = 1,2$ 
  
  \begin{enumerate}
    \item
	What is a probability density function (pdf)? \hfill 1
    \item
	What is P(X=a) in a pdf, where a is an aribitrary number? \hfill 2
    \item  
	Find the marginal probabilities. \hfill 3
    \item
	Find $P(x \vert y), P(x \vert 1)$ and $P(y|4)$ \hfill 4
  \end{enumerate}

 \item
  \textbf{The probability density function of a continuous random variable is}

$$
  f(x) =
\begin{cases}
kx^2+kx+ \frac 18,  & 0 \le x \le 2 \\
0, & otherwise
\end{cases}
$$

  \begin{enumerate}
    \item
	What is a continuous random variable? \hfill 1
    \item
    	Find the value of k \hfill 2
    \item
    	Find the probability that the values of x would lie between1 and 3. \hfill 3
     \item
     	Find the 40th percentile of the distribution and explain.  \hfill 4
  \end{enumerate}
  
   \item
	  \textbf{A professor showed a probability distribution in a class:}
	  
	  \begin{table}[h]
	  \begin{center}
\begin{tabular}{llllll}
x    & 1   & 2 & 3   & 4 & 5   \\ \hline
p(x) & 0.1 & a & 0.3 & b & 0.2 
\end{tabular} \\
\textbf{The value of the arithmetic mean of the distribution is 3.}
\end{center}	
\end{table}


    
  \begin{enumerate}
    \item
	What is the formula of expectation? \hfill 1
    \item
	What is the variance of a constant? Explain logically. \hfill 2
    \item  
	What are the values of a \& b? \hfill 3
    \item
	Find and explain the variance of the distribution. \hfill 4
  \end{enumerate}
  \end{enumerate}

\section{Short Questions}

  \begin{enumerate}
    \item
		What is a continuous random variable?  \hfill 1
    \item
	Question \hfill 1
    \item  
	Question \hfill 1
    \item
	Question \hfill 1
  \end{enumerate}

\chapter{Mathematical Expectation} 
\section{Creative Questions}
  \begin{enumerate}
       \item
	  \textbf{X is a random variable having the below functional form:} 
  
  $P(X) = \frac{6-|7-x|}{k}; x = 1, 2, \cdots,10$ \\
  Y is another variable having the relationship y = 3x+5
  
  \begin{enumerate}
    \item
	What is joint probability? \hfill 1
    \item
	What is the minimum possible value of variance? Why? \hfill 2
    \item  
	Find the value of k. \hfill 3
    \item
	Find E(X) and E(Y). Why are they different? \hfill 4
  \end{enumerate}
  
     \item
	  \textbf{A survey of Television (TV) users at Gulshan in Dhaka was conducted to find how many sets each family use. The following data were obtained:} 
	  
	  	  \begin{table}[h]
	  \begin{center}
\begin{tabular}{llllll}
No of TV set    & 0 & 1  & 2 & 3    \\ \hline
No of family & 10 & 75 & 10 & 5
\end{tabular}
\end{center}	
\end{table}
  
  \begin{enumerate}
    \item
	What is Expectation equivalent to? \hfill 1
    \item
	Can Variance be negative? Why or why not? \hfill 2
    \item  
	Find the variance of the number of TV sets. \hfill 3
    \item
	Find and compoare between arithmetic mean and expectation. \hfill 4
  \end{enumerate}
  \end{enumerate}
\section{Short Questions}

\chapter{Binomial Distribution} 
\section{Creative Questions}
  \begin{enumerate}
 \item
	  \textbf{A farmer plans to store rice seeds for future use. It was found that 8 out of 20 seeds are rotten. He then collected a sample of 15 seeds.} 
  
  \begin{enumerate}
    \item
	What is Bernoulli trial? \hfill 1
    \item
	How are Bernoulli and Binomial distributions related? \hfill 2
    \item  
	What is the probability that at least one seed is rotten out of 15? \hfill 3
    \item
	What is the probability that the number of rotten seeds is greater than the arithmetic mean? \hfill 4
  \end{enumerate}
  
  \end{enumerate}

\section{Short Questions}

\chapter{Poisson Distribution} 
\section{Creative Questions}

  \begin{enumerate}
  \item
  \textbf{In winter, the probability that it rains on a particular day is 0.015. An analyst observes \\ 100 winter days.}
 
  \begin{enumerate}
    \item
	What is an experiment? \hfill 1
    \item
    	When can the Poisson distribution be approximated by the Binomial distribution? \hfill 2
    \item
    	Find, using Binomial distribution, the probability that  it would not rain at all on the \\ observed days. \hfill 3
     \item
     	Find the probability in 3(c) using Poisson distribution.  \hfill 4
  \end{enumerate}

  \item
	  \textbf{BTCL receives 2.5 telephone calls on average from 4 pm to 6 pm. The number of calls received is a random variable. } 
  
  \begin{enumerate}
    \item
	When is Poisson variate applicable? \hfill 1
    \item
	Show conversion criteria and method from Binomial to Poisson distribution. \hfill 2
    \item  
	Find the probability of receiving no more than 3 calls. \hfill 3
    \item
	Find the pattern of calls and show on graph paper.  \hfill 4 \\
	Hint: Find probabilities: P(0), P(1), $\cdots$
  \end{enumerate}

  \end{enumerate}

\section{Short Questions}

\chapter{Normal Distribution} 
\section{Creative Questions}


\section{Short Questions}

\chapter{Index Number} 
\section{Creative Questions}


\section{Short Questions}

\chapter{Sampling} 
\section{Creative Questions}


\section{Short Questions}

\chapter{Vital Statistics} 
\section{Creative Questions}


  \begin{enumerate}
 \item
  \textbf{For projection of population in a future time period, demographers use simple, \\ geometric or exponential growth technique. Each method has its advantages and \\ disadvantages.}

  \begin{enumerate}
    \item
	What is geometric growth? \hfill 1
    \item
    	In geometric growth method, obtain the formula for time required for the population to get \\ doubled [denote rate as r]. \hfill 2
    \item
    	In exponential method, how much unit of time is required for the population to get tripled?  \hfill 3
     \item
     	For projecting (predicting future values), is geometric growth method better than the \\ exponential method? Justify.  \hfill 4
  \end{enumerate}
  
     \item
	  \textbf{Population of Dhaka and Sylhet by different age groups and areas are given below:} 
	  
	  \begin{table}[h]
	  \begin{center}
\begin{tabular}{lllll}
Division & \multicolumn{3}{c}{Age}         & Area ($km^2$) \\
         & 0-14      & 15-64    & 65+      &               \\
Dhaka    & 10,000,00 & 5,00,000 & 5,80,000 & 1,880         \\
Sylhet   & 7,00,000  & 2,70,000 & 4,70,000 & 2,319        
\end{tabular}
\end{center}
\end{table}
  
  \begin{enumerate}
    \item
	Write down the formula of dependency ratio. \hfill 1
    \item
	What is meant by NRR = 0.983? \hfill 2
    \item  
	Find and compare between the dependency ratios of the cities. \hfill 3
    \item
	Based on data, which city is more comfortable for living? \hfill 4
  \end{enumerate}
  \end{enumerate}

\section{Short Questions}

\backmatter
\chapter{Conclusion}
\lipsum[8]

\tableofcontents
\end{document}